% Ad Familiares 4.14 --- headnote
% ad-familiares-04-14, early 46 BC, Rome

\textit{Cicero to Cn. Plancius, written from Rome at the
beginning of 46 BC --- the Perseus dateline reads
\textit{Scr. Romae in. a. 708 (46)}, \textit{in.\ a.}\
marking it as the opening of the year. Plancius is the
same man whom Cicero had defended in 54 in the surviving
\textit{Pro Plancio}; quaestor in Macedonia in 58, he had
been the host of Cicero's exile at Thessalonica, and that
hospitality is the standing debt behind every line of the
correspondence between them. He had taken the Pompeian
side in the civil war and was now living in exile at
Corcyra, where he had sent in the two letters to which
Cicero here replies. The African campaign is still in
prospect; Thapsus is still some months off.}

\textit{The letter turns on a redefinition of \textit{dignitas}.
Plancius has heard that Cicero retains his ``former
standing'' and writes to congratulate him; Cicero begins by
splitting the term. If \textit{dignitas} means thinking
rightly about the state and being approved for it by good
men, he still has it; if it means being able to act on
those thoughts or even defend them in free speech, ``not a
vestige of standing is left to us.'' The second section
turns the consolation outward into a survey of the war's
two possible outcomes, both bleak --- the cruel victory of
angry men on one side, the slaughter of the best citizens
on the other --- and the third pivots to the unspoken
domestic crisis that had occupied Cicero through 47 BC: the
divorce from Terentia and his preparations for the new
marriage to Publilia (``the loyalty of new ones'' against
``the treachery of old ones''). The closing turn back to
Plancius's case is calibrated to his actual situation as a
proscribed Pompeian in Corcyra: Caesar's circle is by now
``reconciled'' to him, the other side ``never was hostile,''
and what Cicero can offer is exertion, counsel, and the
unmistakable signal of continued friendship.}
